\chapter{总结与展望}\label{chap:conclusion}

\section{工作总结}

本文围绕动态立方体抓取任务，构建了一个受控的实验框架，用于系统分析影响模仿学习闭环成功率的关键因素。任务要求机械臂在立方体持续平移和旋转的过程中完成接近、相对位姿对齐、夹爪闭合、保持和搬运——这一过程同时涉及轨迹跟踪、短时运动预测和抓取阶段时机判断。

实验框架的核心设计是目标中心动作块接口：学习策略从历史观测和未来目标条件中预测32步任务空间相对动作块，运行时每16步重规划，单调相位状态机负责阶段推进和夹爪命令，MPC跟踪执行器将任务空间参考转化为关节控制。四类策略——行为克隆（BC）、循环行为克隆（BC-RNN）、动作分块变换器（ACT）和扩散策略（DP）——共享相同的输入输出表示、阶段逻辑和执行器，因此策略之间的比较主要反映动作块生成能力的差异，而非表示或执行器的差异。实验以随机姿态翻滚平移场景（立方体初始姿态均匀随机、绕空间轴旋转）为主要设置，基于5000条专家示教进行训练，以750回合闭环成功率为主要评估标准，报告95\% Wilson置信区间。

\section{主要发现}

在本文的实验设置下，实验结果主要说明了以下三点。

首先，在各项设计因素中，目标中心表示对成功率的影响最为稳定，其影响幅度超过了策略架构差异。以ACT为例：使用目标中心表示时成功率在57\%--59\%之间，切换至世界坐标表示后降至9\%--27\%，差距约为30--50个百分点；相比之下，ACT与BC-RNN在相同表示下的架构差距约为10个百分点。换句话说，在本实验中将全局平移和旋转吸收到相对目标坐标系里，比选择哪种神经网络架构更为根本。需要指出的是，本文专家控制器本身在目标中心坐标系中操作，因此这一发现的普遍程度仍需多专家对比实验来检验。

其次，随机姿态翻滚平移场景需要时序上下文和动作块结构，单步回归无法解决。在本实验的随机姿态翻滚场景中，单步BC几乎完全失败（闭环成功率仅0.1\%），而BC-RNN和ACT分别达到约49\%和59\%。验证实验确认了BC实现本身无误——同一模型在较简单的固定轴自旋场景中可达46\%左右。这说明随机初始姿态和空间轴旋转所要求的抓取面选择、姿态对齐和闭合时机判断，确实超出了单步前馈回归的能力。但时序信息也并非越长越好：BC-RNN在历史长度从1增加到64时出现了显著退化，ACT在历史长度8附近达到峰值。适度的历史窗口和动作块是有效的，但无限制地拉长上下文并不会持续带来收益。

第三，不同类型的控制器对噪声鲁棒性方案的需求不同。基于模型的专家控制器在强噪声条件下完全失效，加入卡尔曼滤波后成功率从0\%恢复到56\%，说明手工规则和显式优化的控制器离不开显式状态估计。而学习策略在相同噪声条件下经卡尔曼滤波后的变化远小于专家：ACT和BC-RNN的变化幅度仅为$\pm$几个百分点，表明这两类时间策略已经通过历史窗口和动作块取得了隐式的时序平滑能力，显式滤波对它们的边际收益有限。在线观测噪声增强训练的效果同样因策略而异——ACT在噪声观测中获得显著改善但无噪性能下降明显，BC-RNN则没有获得可辨识的噪声收益。因此，噪声鲁棒性的实现路径取决于控制器本身的结构特征，不应被视为所有策略共享的通用方案。

综合以上三点，本文实验框架中的动态抓取闭环能力主要取决于表示与时间结构的协同：目标中心表示降低了学习问题的坐标复杂度，时间动作块使随机姿态场景成为可解问题，而噪声鲁棒性则应以匹配控制器类型的方式分别处理。

\section{研究局限与反思}\label{sec:limitations}

回顾整个工作，以下几点限制需要在解读结论时纳入考量。

本文的全部实验均在仿真环境中完成，噪声模型为合成感知扰动，尚未在真实传感器、机械臂接触动力学和通信时延条件下进行验证。任务设置为形状规整的立方体抓取，场景不涉及自然物体、复杂接触和桌面约束，并且抓取点固定（从物体上方抓取），因此结果主要反映轨迹跟踪和时机控制层面的规律，不包含抓取点选择的变化。

更值得注意的一个限制是，所有专家示教均来自同一类控制器——任务空间FSM-MPC——该控制器本身在目标中心坐标系中运行并使用显式阶段结构。这意味着本文关于"表示选择比架构选择更重要"和"时间结构是必需的"等发现，在一定程度上可能反映了策略设计与专家控制范式的一致性，而非动态抓取问题的普遍规律。若采用关节空间强化学习策略、端到端视动策略或人类遥操作示教等不同类型的专家，各设计因素的相对重要性排序可能会有所不同。

此外，几个具体实验设置限制了部分结论的推广范围。学习策略依赖$H_c=50$步的未来目标条件作为输入，在实际部署中预测模块自身的误差会与策略误差叠加，本文未对此进行消融分析。扩散策略在本文中的低成功率仅限于当前低维动作空间和10步DDIM推理设置，不应被推广为对扩散策略的一般性判断。卡尔曼滤波实验中学习策略的增益方向（远小于专家控制器的恢复幅度）一致成立，但具体数值仍需多随机种子的重复实验来确证。

\section{未来工作}

在上述限制的基础上，后续工作可以从以下几个方向推进。首先是硬件验证——将目标中心动作块框架迁移到真实Franka平台，在实际感知误差、控制通信时延和接触动力学的条件下评估各策略的表现差距是否与仿真一致。其次是深化状态估计与学习策略的关系——本文已经区分了两类控制器对显式滤波的不同依赖程度，但二者之间是否存在更优的折中（例如在策略前端仅使用轻量滤波，或让策略直接以滤波状态为输入）仍有待探索。第三是扩展任务分布——引入更复杂的物体形状、桌面接触约束和更宽的目标运动速度范围，检验目标中心表示在更丰富场景中的泛化边界。第四是探索替代结构化执行器的可能性——用学习式低层控制器替换MPC跟踪执行器，或者尝试完全端到端的关节空间策略，从而明确结构化分解（规划/阶段/执行）与端到端学习各自更适合的条件。
